% =========================================================
% CAPITOLO 2 — STATO DELL'ARTE
% =========================================================

\chapter{Stato dell'Arte}
\label{chap:stato_arte}

Nell'ultimo decennio, l'evoluzione delle architetture informatiche ha reso la gestione degli 
standard qualitativi una sfida decisamente più complessa, richiedendo approcci sempre più avanzati.
Con la diffusione dei microservizi, in particolare, garantire le proprietà non funzionali
a runtime è diventato un problema molto più articolato rispetto a quanto avveniva con le 
architetture monolitiche. In un ecosistema dove decine o centinaia di servizi devono collaborare 
per elaborare una singola richiesta, la qualità finale dipende strettamente da come questi 
componenti interagiscono e dalla fluidità del traffico. Ad oggi, le tecniche di rilevamento 
delle anomalie sono efficaci, ma sono caratterizzate da un limite comune: operano sul 
comportamento presente del sistema. Di conseguenza, l'intervallo temporale tra l'origine 
dell'anomalia e la sua rilevazione permette quasi sempre solo un intervento reattivo, 
rendendo difficile un'intervento preventivo. In questo contesto, 
metodologie di assurance e di certificazione forniscono gli strumenti necessari per spostare il focus 
dal semplice rilevamento ad una previsione quantificabile dei malfunzionamenti.

% =========================================================
\section{Sistemi a Microservizi}
\label{sec:microservizi}
% =========================================================

I sistemi distribuiti moderni sono progettati sempre più frequentemente seguendo il modello
a microservizi. Ogni microservizio incapsula un insieme ben definito di funzionalità, comunica 
con il resto del sistema tramite scambio di messaggi e può essere aggiornato o scalato in modo del tutto 
indipendente~\cite{dragoni2017}. Questa scelta progettuale va in contrasto con l'architettura monolitica 
tradizionale, nella quale l'intera logica applicativa risiede in un singolo artefatto, compilato 
e distribuito interamente su ogni istanza. La distinzione tra le due architetture non è solo 
organizzativa; le due architetture sono infatti profondamente differenti nel modo in cui il 
sistema consuma risorse e comunica internamente. Questa distinzione, nelle architetture a 
microservizi, rende molto più complesso garantire e verificare le proprietà non funzionali 
del sistema.

\subsection{Struttura e Modello Comunicativo}
\label{subsec:struttura_comunicazione}

Nelle architetture monolitiche la comunicazione tra sottosistemi avviene attraverso chiamate 
di funzione in memoria. Queste chiamate vengono gestite direttamente dall'hardware locale, 
con latenza nell'ordine dei nanosecondi. Questo meccanismo viene completamente rielaborato 
nel passaggio alle architetture a microservizi, spostando tutto sullo stack di rete. 
In questo caso i servizi interagiscono via \gls{RPC} o \gls{API} RESTful, con tempi di risposta che salgono 
a millisecondi e con maggiori rischi di errori nel passaggio tra un nodo e quello subito successivo. 
Di conseguenza, una singola operazione dell'utente finisce per generare un intero grafo 
di chiamate che coinvolge decine di microservizi diversi. \\
Questa rete di relazioni è rappresentata da un \emph{service dependency graph}, nel quale i vertici 
costituiscono le istanze dei microservizi e gli archi rappresentano le dipendenze comunicative tra di essi~\cite{qiu2020firm}. 
Il comportamento che si manifesta a runtime su questo grafo è dinamico. Questa struttura, infatti, 
evolve in risposta a deploy automatici, scaling orizzontale e configurazioni di routing che variano 
al variare del carico: nuove istanze di un servizio vengono aggiunte e rimosse, i load balancer redirigono il traffico, 
i circuit breaker aprono percorsi alternativi quando un servizio è in stato degradato. 
Questa dinamicità, che nei microservizi si manifesta a una velocità e a una granularità 
senza precedenti rispetto alle architetture \gls{SOA} tradizionali, costringe i sistemi di monitoraggio a 
trattare la topologia come un'informazione strutturale viva, e non come una semplice mappa statica.

La granularità fine dei microservizi, pur facilitando la manutenzione, introduce un overhead 
di rete senza equivalenti nell'architettura monolitica; difatti, misurazioni empiriche 
sull'applicazione Social Network di DeathStarBench~\cite{gan2019deathstar} evidenziano come la gestione delle 
comunicazioni e l'elaborazione delle richieste di rete possano costituire il 36,3\% del tempo totale 
di esecuzione end-to-end. E' quindi possibile affermare che la qualità del servizio end-to-end 
non è determinata solo dalla capacità computazionale di ciascun nodo preso singolarmente; 
bensì dipende anche in modo rilevante dalla qualità di ogni singola interazione tra 
coppie di nodi. In questo ambito, metriche come latenza e tasso di errore sugli archi diventano 
indicatori primari per la compliance del sistema.
Per ottenere una visibilità strutturata sulle interazioni tra microservizi si ricorre 
al distributed tracing. Strumenti come Jaeger\footnote{\url{https://www.jaegertracing.io/}} e Zipkin\footnote{\url{https://zipkin.io/}}, 
ormai standard in ambito cloud-native, 
permettono infatti di ricostruire l'intero flusso di esecuzione di una richiesta attraverso 
un grafo di span organizzati gerarchicamente. Questi span vengono annotati con timestamp 
di inizio e fine dell'interazione, identificatori di servizio e metadati di esito. L'integrazione di 
questi dati con le metriche di sistema raccolte da agenti come Prometheus\footnote{\url{https://prometheus.io/}}, che campionano 
continuamente l'utilizzo delle risorse a livello di container, permette di ottenere una 
rappresentazione del comportamento del sistema che è contemporaneamente fine-grained a livello 
di interazione e persistente nel tempo come serie storiche. Per unificare queste due dimensioni 
nasce il formalismo degli \gls{ATG} adottato in questo framework, caratterizzato sia dai dati puntuali dei 
singoli componenti sia dalla struttura temporale delle loro interazioni. In questo contesto, 
un esempio concreto di adozione del distributed tracing è DeathStarBench~\cite{gan2019deathstar}. Quest'ultimo adotta un 
sistema di tracciamento distribuito basato su Jaeger tramite le \gls{API} OpenTracing\footnote{\url{https://opentracing.io/} - 
La specifica OpenTracing è stata archiviata nel 2022 e incorporata in OpenTelemetry.}, capace di tracciare 
le richieste a granularità \gls{RPC}, associare le chiamate appartenenti alla stessa richiesta end-to-end 
e registrare le tracce in un database centralizzato. 

\subsection{Proprietà Non Funzionali nei Sistemi a Microservizi}
\label{subsec:proprieta_nonfunzionali}

Le proprietà non funzionali di maggiore rilevanza per il problema del forecasting affrontato 
in questa tesi sono la latenza end-to-end, la disponibilità dei percorsi di servizio e il 
throughput. Queste caratteristiche rappresentano infatti dimensioni critiche sia nella ricerca 
accademica che nell'operatività industriale~\cite{auer2021}. Ciascuna di queste proprietà presenta 
caratteristiche che la distinguono dalla corrispondente proprietà in architettura monolitica.

La \textbf{latenza end-to-end} di una richiesta è determinata dal percorso critico, ossia
il cammino di durata massima dalla ricezione della richiesta alla produzione della risposta.
Nei microservizi, il percorso critico non corrisponde necessariamente al percorso più lungo
nel service dependency graph; il percorso critico, infatti, cambia continuamente in base al
carico istantaneo sui singoli nodi e alle fluttuazioni della rete. Una stessa richiesta può
seguire percorsi critici differenti in momenti diversi, a seconda di quale componente stia
subendo un rallentamento in quel preciso istante~\cite{qiu2020firm}. La latenza, di fatto,
smette di essere un valore stabile e diventa una variabile stocastica che evolve nel tempo.

Un discorso analogo vale per la \textbf{disponibilità}. Se un percorso di servizio attraversa
$n$ microservizi in serie, la sua disponibilità complessiva è il prodotto delle disponibilità
dei singoli componenti. Questo ha un impatto pratico enorme:
con venti microservizi, ciascuno con una disponibilità individuale del 99,9\%, la disponibilità
totale del percorso crolla al 98\%. Un valore del genere è spesso inaccettabile per sistemi
critici~\cite{gan2019deathstar}. Questo fenomeno giustifica l'adozione di pattern
architetturali di resilienza come il circuit breaker, per deviare il traffico dai nodi guasti,
e il bulkhead, che isola i guasti impedendone la propagazione~\cite{qiu2020firm}. Questi
meccanismi, però, modificano la topologia del sistema a runtime, rendendo ancora più complessa
la capacità di monitorarne l'evoluzione strutturale.

Infine, il \textbf{\textit{throughput}} del sistema è limitato dall'elemento che costituisce il
``collo di bottiglia''. Il servizio con la capacità di elaborazione più bassa vincola la
portata dell'intero percorso. La situazione si complica ulteriormente quando più flussi
competono per la stessa risorsa, come un database condiviso. In questi casi, la contesa riduce
il throughput disponibile per ciascun flusso in modo non lineare. Senza una modellazione
esplicita di queste dipendenze, prevedere il throughput reale diventa estremamente difficile.

\subsection{Propagazione a Cascata e Sfide per il Monitoraggio}
\label{subsec:cascading}

Uno dei fenomeni più critici nei sistemi a microservizi è la propagazione a cascata delle 
violazioni di \gls{QoS}. Un singolo microservizio che accumula ritardi per saturazione delle risorse, 
a causa di picchi di traffico o contesa su una risorsa condivisa, fa sì che i componenti che 
dipendono da esso vedano un aumento della propria latenza di risposta. Quando la latenza va al 
di sopra dei limiti di timeout configurati, i servizi chiamanti spesso tentano nuovamente l'operazione; 
tuttavia, ciò amplifica ulteriormente il carico sul servizio già degradato, causando una reazione a 
catena che risale le dipendenze coinvolgendo anche i servizi originariamente sani. 
La gestione non ottimale di una singola dipendenza tra microservizi è in grado di degradare 
la latenza di un'intera applicazione~\cite{gan2019deathstar}. In questo caso, anche dopo aver rimosso la causa del problema, 
il sistema richiede spesso periodi prolungati per tornare a uno stato di equilibrio.

Questo meccanismo di propagazione rende inadeguati i sistemi di monitoraggio attuali, che si
limitano ad osservare componenti individuali in isolamento. Quando viene rilevata un'anomalia
tra due servizi del service dependency graph, la causa radice può trovarsi in un nodo $w$
topologicamente distante, oppure in un flusso applicativo del tutto separato che condivide
con il percorso monitorato una risorsa critica, come un database o un message broker.
Identificare l'origine di un guasto richiede quindi la capacità di ragionare sulla struttura
complessiva del grafo e di seguire il percorso di propagazione degli effetti, superando la
logica del semplice rilevamento di un'anomalia al tempo $t$.
Un'ulteriore complicazione deriva dall'eterogeneità tecnologica dei microservizi. 
Poichè ciascun servizio può essere sviluppato con linguaggi e runtime differenti, 
ognuno di essi reagisce al carico in modo proprio, con profili di hotspot e latenza che 
variano a seconda dell'ambiente di esecuzione~\cite{gan2019deathstar}. Tutto ciò ha una conseguenza: 
le metriche che segnalano un problema imminente non sono universali. 
Un indicatore fondamentale per un certo componente può risultare del tutto irrilevante per un altro. 
Un sistema di monitoraggio efficace deve quindi incorporare questa varietà, evitando di appiattire 
tutto su soglie di allarme generiche e uniformi. 

Dal punto di vista del forecasting, la propagazione a cascata ha un'implicazione che va oltre 
la semplice diagnosi; esiste infatti un ritardo sistematico tra l'insorgere della causa radice 
e il momento in cui la violazione diventa visibile esternamente nelle metriche di compliance. 
Quando un componente inizia a degenerare internamente, le metriche osservate a
valle possono apparire normali per un tempo prolungato, nascondendo il degrado che sta
crescendo internamente al nodo. Questa latenza è, allo stesso tempo, un ostacolo e
un'opportunità. È un ostacolo perché rende il monitoraggio reattivo ``cieco'' alla causa in
formazione; è però un'opportunità perché quella stessa finestra temporale, se sfruttata
attraverso il monitoraggio delle metriche interne di nodo, rappresenta il \emph{lead time}
necessario per un intervento preventivo.
Tutto questo chiarisce perché, per il forecasting della non-compliance, non basti analizzare
serie temporali isolate. È necessaria, infatti, una visione d'insieme che solo una
rappresentazione a grafo con feature eterogenee, su nodi e archi, può offrire. Il
Capitolo~\ref{chap:metodologia} formalizza questa necessità introducendo gli \gls{ATG} e la
loro integrazione con la struttura semantica degli ipergrafi.

% =========================================================
\section{Assessment dei Sistemi Distribuiti}
\label{sec:assessment}
% =========================================================

Nell'ambito dei sistemi distribuiti, il termine assessment rappresenta l'insieme delle attività 
di osservazione e analisi del comportamento di un sistema a runtime. Lo scopo di queste attività 
è quello di individuare eventuali anomalie e verificare la conformità del sistema stesso rispetto 
a specifici requisiti attesi. L'assessment si articola su tre diversi livelli di astrazione crescente: 
il monitoraggio passivo, il rilevamento delle anomalie e l'analisi delle cause radice. 
Ciascun livello si fonda su quello precedente, aggiungendo capacità diagnostiche che i livelli 
inferiori non sono in grado di fornire autonomamente.
Nel corso degli anni, la letteratura in questo campo ha prodotto una vasta gamma di soluzioni, 
raggruppabili in quattro categorie principali: metodi statistici classici, metodi basati su \gls{ML}, 
metodi fondati su grafi e \gls{GNN}, e tecniche di root cause analysis. Una direzione più recente, 
particolarmente rilevante per il problema del forecasting affrontato in questa tesi, 
è quella dei grafi temporali con feature esplicite sui nodi, discussa nella sezione~\ref{subsec:atg_background}.

\subsection{Monitoraggio Passivo e Osservabilità}
\label{subsec:monitoring}

Il monitoraggio passivo rappresenta la base dell'osservabilità di un sistema. 
L'utilizzo dell'infrastruttura di telemetria permette infatti di raccogliere in modo costante 
le metriche di sistema, come l'utilizzo della \gls{CPU} o la latenza, rendendole disponibili per le analisi successive. 
Tale pratica è ormai standardizzata grazie a soluzioni consolidate come Prometheus per la gestione delle metriche, 
OpenTelemetry\footnote{\url{https://opentelemetry.io/}} per l'instrumentazione unificata e Jaeger per il tracciamento delle richieste distribuite. 
Tuttavia, questi strumenti soffrono di un problema comune: la quantità di dati generati supera 
di gran lunga quella dei sistemi tradizionali. Un'architettura a microservizi di medie dimensioni 
può arrivare infatti a produrre centinaia di serie temporali e milioni di span di traccia 
per ogni ora di operatività~\cite{gan2019deathstar}. Questa quantità di dati genera un problema di selezione. 
Non tutte le metriche sono ugualmente rilevanti per le diverse proprietà da garantire e 
l'analisi indiscriminata di ogni serie disponibile non è computazionalmente praticabile in tempo reale.

Il monitoraggio passivo non genera diagnosi in modo autonomo, ma si limita a fornire la base
informativa per i livelli analitici superiori. Questi ultimi devono essere progettati
specificamente per identificare i dati utili, separandoli dal rumore di fondo della telemetria.
La scelta di quali metriche osservare, della frequenza di campionamento e del livello di
dettaglio costituisce a tutti gli effetti un problema di progettazione che ha conseguenze
sull'efficacia delle analisi a valle.

\subsection{Rilevamento delle Anomalie}
\label{subsec:anomaly_detection}

Il rilevamento delle anomalie funge da filtro critico che trasforma la grande quantità di 
serie temporali grezze in precisi segnali di allerta, identificando quindi i momenti in cui 
il comportamento del sistema devia dalla normale operatività. Le tecniche disponibili spaziano 
a partire da approcci statistici tradizionali, come z-score o CUSUM, fino a modelli di apprendimento 
automatico in grado di processare pattern multidimensionali che 
l'analisi univariata non sarebbe in grado di intercettare.

I \textbf{metodi statistici classici} permettono di garantire un livello di interpretabilità elevato 
e una calibrazione semplice. Fissare una soglia di allarme sulla latenza di un arco, 
ad esempio, rappresenta il modo più semplice per mappare un vincolo contrattuale di \gls{SLA} 
in una semplice operazione matematica di confronto. Tuttavia, questi strumenti tradizionali 
sono caratterizzati da una forte rigidità nei confronti di fluttuazioni fisiologiche 
del comportamento dei sistemi in analisi. Un picco di latenza durante le ore di punta potrebbe 
essere del tutto fisiologico, mentre lo stesso valore registrato subito dopo un deploy 
assumerebbe le sembianze di una criticità. Seppur vi siano metodi statistici più elastici ed adattivi, 
come il CUSUM, che permettono di accumulare deviazioni dalla baseline per rilevare shift persistenti 
anziché picchi isolati, il problema di fondo persiste. Difatti, pur essendo più elastici ed adattivi, 
si tratta comunque di metodi univariati, incapaci quindi di tenere conto delle dipendenze 
tra metriche di componenti distinti.

Per superare i limiti dell'analisi univariata si utilizzano i \textbf{metodi basati su machine learning}. 
Questi modelli non si limitano ad osservare le singole metriche, bensì apprendono il comportamento 
fisiologico del sistema, sotto forma di struttura multidimensionale, a partire dai dati storici. 
Algoritmi come l'Isolation Forest, ad esempio, sono in grado di isolare i punti anomali nello spazio delle feature, 
procedendo attraverso partizionamenti ricorsivi. Questa logica permette di individuare combinazioni 
di valori che presi singolarmente sembrerebbero normali, ma che nel loro insieme rivelano uno stato di allerta.

Un'ulteriore evoluzione è rappresentata dai metodi basati su \textbf{\gls{GNN}}. Questi metodi estendono le capacità dei metodi 
precedenti modellando le interazioni tra microservizi attraverso la struttura del grafo di dipendenze. 
Operando su questa rappresentazione, è possibile rilevare anomalie che si manifestano come pattern 
insoliti tra componenti, non solo come deviazioni di singole metriche. Valutazioni sperimentali su 
dataset con bottleneck multipli mostrano miglioramenti dell'F1 score superiori al 40\% rispetto ad 
approcci che non sfruttano reti neurali su grafi per modellare le interazioni strutturali tra microservizi~\cite{somashekar2024gamma}. 

\textbf{GAMMA (Somashekar et al., 2024).} Il framework GAMMA affronta il rilevamento di bottleneck multipli simultanei nelle architetture 
a microservizi, uno scenario spesso trascurato dalla letteratura precedente che tende a concentrarsi 
sul singolo bottleneck. L'architettura sfrutta una \gls{GAT} che apprende le interazioni tra microservizi 
a partire dalle tracce distribuite e dalle metriche di sistema, integrandola con un \emph{mixture of experts} 
che assegna a ogni microservizio un classificatore specializzato nel riconoscimento di specifici profili di degrado. 
L'addestramento viene eseguito in modo congiunto sui task di \emph{anomaly detection} e \emph{bottleneck localization} 
attraverso una funzione di loss unificata. Questo approccio permette alle rappresentazioni apprese nei due 
sotto compiti di rinforzarsi reciprocamente durante il processo. Le valutazioni condotte su DeathStarBench~\cite{gan2019deathstar} 
evidenziano un F1 score fino a 0,92 per il rilevamento e fino a 0,89 per la localizzazione, 
con un incremento del 46\% nelle prestazioni di localizzazione rispetto a FIRM in presenza di bottleneck multipli. 
Il limite strutturale di GAMMA, condiviso con tutti i metodi di questa categoria, è la natura esclusivamente reattiva. 
Il sistema opera infatti solo sul comportamento presente del sistema, tale per cui la classificazione dello 
stato avviene al tempo $t$, senza produrre stime sull'evoluzione futura dei segnali nel tempo. Di conseguenza, 
il \emph{lead time} per un intervento preventivo risulta nullo per costruzione.

\textbf{Anomaly Detection and Diagnosis (Du et al., 2018).} L'approccio~\cite{du2018} introduce un 
sistema di \emph{anomaly detection} per microservizi containerizzati che affronta un problema di osservabilità pratico: 
quali metriche di sistema raccogliere a livello di container e come valutare automaticamente se il comportamento 
è anomalo senza conoscere a priori la struttura interna dei servizi. Il sistema sfrutta agenti di 
monitoraggio per raccogliere dati di performance, come \gls{CPU} e rete, e utilizza classificatori 
supervisionati come Random Forest, \gls{kNN} e \gls{SVM}. Per addestrare i modelli a distinguere e classificare le anomalie, 
il sistema si affida a dataset generati tramite fault injection controllata. 
Nelle valutazioni sperimentali sul dataset B, Random Forest e kNN raggiungono precision e recall superiori a 0,97. 
Sul dataset aggregato C, che combina le due sessioni sperimentali, i valori scendono a circa 0,96 per kNN e 0,93/0,91 
per Random Forest; nella validazione per singolo servizio i risultati calano ulteriormente. 
Il limite principale di questa metodologia risiede nella stretta dipendenza da dati etichettati prodotti artificialmente. 
In scenari di produzione reali, caratterizzati da distribuzioni di guasto variabili o da eventi mai 
osservati durante il training, la capacità di generalizzazione risulta sensibilmente ridotta rispetto a soluzioni 
\textit{unsupervised} o \textit{semi-supervised}.

\subsection{Analisi delle Cause Radice}
\label{subsec:rca}

L'analisi delle cause radice estende il rilevamento
identificando non solo che un'anomalia è in corso, ma anche quale componente ne è la causa
originaria. La distinzione è operativamente critica: un operatore che riceve l'alert ``latenza
elevata sull'arco $e_2$'' non ha informazioni sufficienti per decidere se intervenire sul
servizio sorgente, sul servizio destinazione o su una risorsa condivisa che entrambi contendono.
Lo scopo della \gls{RCA} è produrre un'indicazione azionabile che consenta di dirigere
l'intervento correttivo verso la causa vera, non verso il sintomo visibile.
Le tecniche attuali di \gls{RCA} nei microservizi combinano l'analisi della topologia del
grafo di dipendenze con metodi statistici di causalità, consentendo di risalire dal sintomo
osservato alla componente responsabile. Sistemi che integrano analisi del percorso critico con
classificatori supervisionati per la localizzazione della responsabilità \gls{SLO} dimostrano
che è possibile identificare il microservizio responsabile di una violazione con accuracy media
dell'ordine del 93\% su benchmark standardizzati~\cite{qiu2020firm}.

Resta però un limite: anche le tecniche di \gls{RCA} più avanzate operano su ciò che è già
accaduto. La diagnosi riguarda un'anomalia presente o recente; non permette di predire
un'anomalia futura. In altri termini, l'output prodotto rappresenta una fotografia del guasto
e non un forecast della sua possibile insorgenza futura.
Un limite ulteriore, che si rivela particolarmente rilevante per il problema della tesi,
riguarda le cause radice che si propagano tra flussi applicativi distinti. 
Quando due percorsi certificati condividono una risorsa, come un database o un message
broker, il degrado prodotto dal carico di un flusso può manifestarsi come violazione
di latenza nell'altro flusso. Nessuno dei due sistemi di monitoraggio, analizzato
individualmente, è in grado di ricostruire questa catena causale. Il sistema
di monitoraggio del flusso osserva un incremento della latenza sul proprio percorso e
può correttamente attribuirlo a una saturazione del database condiviso; ma per capire perché
il database è saturo occorre correlare con le metriche del flusso appartenente a un perimetro di compliance diverso. 
In assenza di una rappresentazione esplicita
delle dipendenze tra proprietà certificate diverse, questa correlazione rimane implicita e la
diagnosi si ferma al sintomo, non alla causa strutturale. Il framework proposto affronta
questo problema non come un'estensione della \gls{RCA}, ma incorporando nella struttura del modello 
la conoscenza topologica delle risorse condivise; cosi facendo l'interferenza diventa visibile prima ancora
che si traduca in anomalie osservabili sui dati.

\textbf{FIRM (Qiu et al., 2020).} FIRM~\cite{qiu2020firm} propone un framework di 
localizzazione delle violazioni \gls{SLO} e mitigazione delle risorse orientato ai 
microservizi, basato sull'integrazione di componenti complementari. 
Il processo di localizzazione si basa sul concetto di percorso critico dinamico all'interno 
del grafo di esecuzione di ogni richiesta. Nello specifico, un classificatore \gls{SVM} analizza le 
metriche di variabilità del \emph{sojourn time}, ossia il tempo di permanenza della richiesta nel servizio, 
lungo tale percorso per identificare l'istanza di microservizio responsabile del rallentamento. 
La mitigazione utilizza un agente di \gls{DDPG}, addestrato con transfer learning tra configurazioni 
diverse di workload, per modulare in tempo reale l'allocazione delle risorse condivise 
sul microservizio localizzato. Su quattro benchmark di microservizi, FIRM riduce le violazioni 
\gls{SLO} fino a 16 volte rispetto al meccanismo di autoscaling nativo di Kubernetes\footnote{\url{https://kubernetes.io/}}; 
inoltre, localizza la causa radice della violazione con accuracy media del 93\%~\cite{qiu2020firm}. 
Tuttavia, nel contesto di questa tesi, il limite di FIRM è duplice. Il primo riguarda 
l'assunzione implicita di single bottleneck, tale per cui un solo microservizio è 
responsabile della violazione in un dato istante; mentre il secondo riguarda l'assenza 
di capacità predittive, in quanto il sistema interviene solo a violazione avvenuta o in corso.

\textbf{Root Cause Discovery (NeurIPS 2022).} L'algoritmo di \emph{causal discovery}~\cite{ikram2022rcd} denominato \gls{RCD}, 
introduce due innovazioni chiave per rendere scalabile la \gls{RCA} all'interno dei sistemi a microservizi. 
La prima è di ordine concettuale, il guasto viene trattato infatti come un intervento sul nodo di origine 
all'interno del grafo causale; ciò permette di distinguere tra i dati raccolti durante 
il malfunzionamento, ossia l'\emph{interventional dataset}, e quelli relativi al normale funzionamento, 
ossia l'\emph{observational dataset}. Il successivo confronto tra le distribuzioni dei due dataset permette 
di ottenere informazioni utili sulla causa radice che i soli dati osservativi non potrebbero fornire. 
La seconda innovazione riguarda l'efficienza algoritmica. Invece di mappare 
l'intera struttura causale del sistema, \gls{RCD} esplora in modo gerarchico e localizzato solo le porzioni 
del grafo statisticamente connesse al nodo in errore. Questo approccio riduce sensibilmente il numero 
di test di indipendenza condizionale necessari. I test condotti su sistemi da 500 nodi mostrano tempi 
di esecuzione estremamente ridotti: circa 22 secondi per identificare le prime cinque candidate, 
contro gli oltre 150 minuti richiesti dall'algoritmo $\Psi$-PC~\cite{ikram2022rcd}. L'importanza di tali strumenti emerge chiaramente: 
senza l'utilizzo di essi, l'identificazione della causa radice richiede in media 3 ore. 
Tuttavia, \gls{RCD}, come tutti i metodi di \gls{RCA}, permette di identificare la causa di un guasto che si è già verificato, 
ma non produce alcuna stima della probabilità o del tempo di occorrenza di guasti futuri.

\subsection{Verso la Modellazione Temporale con Feature di Nodo}
\label{subsec:atg_background}

Tutte le limitazioni evidenziate finora conducono verso una chiara ed unica necessità, ossia superare 
la logica dell'istante $t$ per la modellazione dell'evoluzione temporale a favore di un modello che integri 
lo stato interno dei nodi con la struttura delle loro interazioni. Nonostante la ricerca sui grafi dinamici 
affronti questa problematica, nella sua formulazione classica tende comunque a concentrare l'informazione 
primaria sulle interazioni. In questo contesto, un grafo temporale viene tradizionalmente formalizzato come 
una sequenza di eventi $(u, v, t)$, dove l'esistenza stessa del sistema è definita dagli scambi di messaggi 
tra entità, e lo stato di un nodo è derivato esclusivamente dalla sua storia di interazioni. 
Un nodo assume infatti rilevanza solo in funzione di ciò che invia o riceve sugli archi.
Questo approccio è sufficiente per studi topologici, ma risulta limitante per il monitoraggio
predittivo di sistemi distribuiti reali, nei quali:

\begin{itemize}
  \item i componenti possiedono uno stato interno persistente (utilizzo della \gls{CPU},
        occupazione di memoria, saturazione del connection pool) che evolve in modo
        indipendente dal traffico in transito;
  \item il degrado nasce spesso prima che si manifesti sulle interazioni: un componente
        può degradarsi internamente senza che le metriche di arco abbiano ancora superato
        le soglie di allerta;
  \item la violazione dei requisiti di compliance è solitamente preceduta da una deriva
        comportamentale osservabile sullo stato dei nodi e non sugli archi.
\end{itemize}

La ricerca recente nell'ambito del graph representation learning ha introdotto grafi temporali 
dotati di feature esplicite sui nodi~\cite{rossi2020}, estendendo così la formulazione classica. 
Nelle \gls{TGN} viene infatti formalizzato il concetto di \textbf{Node Memory}, che permette a ogni nodo di 
mantenere uno stato latente evolutivo. Questo stato può essere aggiornato attraverso 
due tipologie di eventi, ossia le interazioni tra nodi sugli archi e gli eventi locali al nodo, 
indipendenti dal traffico osservato. Gli eventi locali al nodo sono fondamentali per i sistemi distribuiti; 
difatti, le metriche interne di un microservizio, come l'occupazione dell'heap o la saturazione delle code, 
possono mostrare segni di deterioramento anche in assenza di chiamate di servizio attive. \\
L'efficacia di questo approccio trova conferma in contesti analoghi, come lo studio delle reti collaborative. 
Nei dataset di Wikipedia utilizzati per valutare modelli come JODIE o \gls{TGN}, le modifiche alle 
pagine sono modellate come eventi temporali sugli archi, ad esempio l'interazione ``utente modifica pagina''. 
In questo caso, gli utenti, rappresentati dai nodi, conservano uno stato latente di affidabilità 
che può degradare nel tempo; difatti, quando un utente effettua modifiche troppo rapide o incoerenti, 
il suo stato interno degrada verso una condizione di sospetto vandalismo prima ancora che venga 
riscontrata una violazione esplicita delle policy. Il parallelismo con i sistemi distribuiti in questo caso è diretto. 
Un microservizio può accumulare segnali interni di degrado, come memory leak o saturazione delle risorse, 
indipendentemente dalle chiamate \gls{API} in corso, anticipando il momento in cui le interazioni esterne inizieranno a deteriorarsi.

Un'ulteriore conferma della validità di questo approccio arriva dai sistemi di previsione del 
traffico su reti stradali. In questi modelli, i nodi rappresentano incroci fisici caratterizzati 
da uno stato evolutivo proprio, come la densità veicolare o la velocità media; 
le strade rappresentano invece gli archi di collegamento. Spesso un incrocio raggiunge uno stato critico di congestione
prima che il rallentamento si propaghi visibilmente ai tratti stradali adiacenti; il fenomeno
nasce quindi come evento locale al nodo, non ancora osservabile sugli archi in quel preciso
istante. Nei sistemi distribuiti si osserva un comportamento simile, ad esempio in un gateway
che accumula una coda crescente di richieste. Questo accumulo funge da precursore per i timeout
futuri, ma inizialmente resta un segnale di degrado confinato interamente nello stato del nodo,
senza influenzare in modo immediato la latenza misurata sugli archi incidenti. \\
Tali evidenze confermano che la modellazione esplicita dello stato interno dei componenti, e
non solo delle loro interazioni, costituisce un requisito strutturale per qualsiasi forma di
monitoraggio predittivo. La distinzione tra feature di nodo, che persistono anche in assenza
di traffico, e feature di arco, effimere e dipendenti dal contesto comunicativo, rappresenta
il nucleo della struttura duale adottata dal framework presentato nel Capitolo~\ref{chap:metodologia}. 

% \subsection{Il Limite Fondamentale: Reattività senza Anticipo}
% \label{subsec:lead_time_zero}

% Monitoraggio, anomaly detection e \gls{RCA} condividono un limite strutturale che questa tesi
% si propone di superare: operano esclusivamente sull'istante $t$, anziché su un orizzonte
% temporale. Un sistema di anomaly detection si limita a segnalare che il comportamento al tempo
% $t$ è anomalo, mentre la \gls{RCA} identifica la causa di un evento già manifesto. In questi casi, 
% il preavviso effettivo si riduce al puro ritardo di 
% rilevamento tra il momento di origine dell'anomalia e la sua 
% identificazione, tipicamente dell'ordine di secondi o di pochi minuti.
% Questo preavviso risulta spesso insufficiente per attivare strategie di mitigazione efficaci.
% Operazioni come il provisioning di risorse aggiuntive, il reindirizzamento del traffico verso
% repliche sane o la limitazione proattiva del carico su percorsi a rischio di congestione
% richiedono infatti tempi di reazione superiori al margine offerto da un rilevamento puramente
% reattivo. \\
% FIRM riconosce esplicitamente che il forecasting delle violazioni  rimane un problema aperto, identificandolo
% come direzione per lavori futuri~\cite{qiu2020firm}. GAMMA, pur migliorando le capacità di
% detection e localizzazione delle cause in scenari multi-bottleneck, opera anch'esso
% esclusivamente sul comportamento presente, senza quindi produrre stime sull'evoluzione
% temporale del sistema~\cite{somashekar2024gamma}.\\
% Il forecasting della non-compliance che questa tesi propone colma precisamente questo gap.
% Invece di classificare lo stato corrente come normale o anomalo, il framework stima l'evoluzione
% futura delle metriche di compliance, calcolando il tempo residuo prima del superamento di una
% soglia certificata. Il \emph{lead time} generato non è quindi un semplice residuo del ritardo
% di rilevamento, ma una stima quantitativa dell'orizzonte disponibile per l'intervento preventivo.

% =========================================================
\section{Assurance e Certificazione di Proprietà Non Funzionali}
\label{sec:certificazione}
% =========================================================

Il termine \emph{assurance} designa l'insieme delle pratiche consolidate per dimostrare che
un sistema si comporti correttamente rispetto a un insieme di requisiti
non funzionali~\cite{ardagna2023manifesto}. A differenza dell'assessment, che si limita a osservare
il comportamento del sistema a runtime per rilevarne le deviazioni, l'assurance produce
garanzie basate su evidenze raccolte secondo un processo strutturato. Tale percorso è governato
da standard documentati e risulta verificabile da terze parti. La certificazione rappresenta
la forma più matura di assurance. Essa produce un artefatto formale, il certificato, che
attesta la conformità di un sistema a una specifica proprietà non funzionale. Questo avviene
seguendo un iter definito da una \gls{CA} e verificato da laboratori
accreditati~\cite{anisetti2023multidim}.

\subsection{Elementi del Processo di Certificazione}
\label{subsec:processo_certificazione}

Un processo di certificazione si basa su tre costrutti fondamentali~\cite{anisetti2023continuous}:

\begin{itemize}
  \item \textbf{Proprietà non funzionale}: modella il requisito da verificare, specificando 
      il nome della proprietà, ad esempio \emph{performance}, e l'insieme degli attributi che definiscono 
      le condizioni di soddisfacimento, come la \emph{latenza massima};
  \item \textbf{\gls{ToC}}: rappresenta il sistema da certificare. 
      Quest'ultimo è modellato come un insieme di componenti raggruppati secondo la propria funzionalità;
  \item \textbf{\gls{CM}}: specifica le attività da eseguire sul \gls{ToC} per raccogliere le
        evidenze necessarie a dimostrare che la proprietà $p$ sia effettivamente soddisfatta.
\end{itemize}

Il \gls{CM}, una volta definito dalla \gls{CA}, viene istanziato su delega da un laboratorio 
accreditato. Quest'ultimo ha il compito di raccogliere le evidenze secondo le attività 
specificate nel \gls{CM}. La raccolta avviene attraverso un \gls{ECM}, che si occupa di disciplinare 
la modalità, i criteri e gli strumenti di raccolta dal \gls{ToC}. Il processo è
illustrato nella Figura~\ref{fig:certificate}.\\
Non appena il volume di evidenze raccolto risulta sufficiente per i criteri del CM, la CA rilascia un certificato C. 
Tale documento include il modello di certificazione adottato, le evidenze raccolte e garantisce la catena 
di fiducia attraverso una firma digitale. La validità del certificato rilasciato dipende dalla correttezza 
di tre assunzioni~\cite{anisetti2023continuous}: il \gls{CM} deve essere prodotto da una \gls{CA} fidata, 
il \gls{CM} deve rappresentare fedelmente il \gls{ToC} e la proprietà certificata ed infine le evidenze raccolte devono essere sufficienti 
a dimostrare la proprietà. Il certificato rimane valido finché queste tre condizioni continuano ad essere soddisfatte.

\begin{figure}[H]
	\centering
	\includegraphics[width = 0.60\textwidth, height=0.45\textheight, trim={0 10mm 0 0}, clip]{images/certificate}
	\caption[Modello di certificazione]{Modello di certificazione~\cite{anisetti2023continuous}}
	\label{fig:certificate}
\end{figure}
\begin{figure}[H]
	\centering
	\includegraphics[width = 0.8\textwidth, height=0.35\textheight,trim={0 0 0 0}, clip]{images/certificate_usage}
	\caption[Utilizzo della certificazione]{Utilizzo della certificazione~\cite{anisetti2023multidim}}
	\label{fig:certificate_usage}
\end{figure}

Schemi di questo tipo trovano applicazione consolidata nella verifica di proprietà di
confidenzialità, integrità e disponibilità in ambienti cloud, oltre che per la selezione di
servizi basata su \gls{SLA}~\cite{anisetti2023multidim}. La Figura~\ref{fig:certificate_usage} 
illustra come, nel contesto dei sistemi distribuiti moderni, le certificazioni vengano utilizzate. 
Il Service Provider implementa il servizio, la CA esegue il Certification Model contro di esso, 
e il certificato risultante viene depositato in un repository condiviso. L'End User 
consulta questo repository per selezionare il servizio più idoneo ai propri requisiti, 
sfruttando le proprietà certificate come criteri di confronto oggettivi. Questo modello funziona 
però a condizione che il sistema certificato rimanga sostanzialmente stabile nel tempo, un'assunzione 
che i sistemi a microservizi mettono in crisi.

\subsection{Limiti degli Schemi Tradizionali nei Sistemi Distribuiti Dinamici}
\label{subsec:limiti_tradizionali}

Gli schemi di certificazione tradizionali si basano su assunzioni che i sistemi distribuiti 
moderni violano sistematicamente. La più critica è la staticità del \gls{ToC}, tale per cui il \gls{CM} 
descrive il sistema nella configurazione osservata al momento della certificazione, assumendo 
che tale configurazione persista nel tempo~\cite{ardagna2023manifesto}. Nelle architetture a microservizi questa 
assunzione non regge; difatti, fenomeni come deploy automatici, scaling orizzontale, 
failover e riconfigurazioni orchestrate, modificano costantemente l'insieme dei componenti attivi, 
la loro topologia di comunicazione e le prestazioni complessive. Un certificato rilasciato per 
una specifica configurazione architetturale non è quindi automaticamente valido per la 
configurazione successiva, anche quando le proprietà funzionali del sistema rimangono invariate. 

La seconda limitazione riguarda il meccanismo di innesco della ricertificazione. Gli schemi
attuali reagiscono principalmente a modifiche del codice sorgente o a scadenze temporali
predefinite, ignorando le variazioni comportamentali del sistema non relative a modifiche
strutturali esplicite~\cite{anisetti2023continuous}. Un sistema che subisce un degrado graduale
delle proprie performance, causato ad esempio dalla saturazione progressiva delle risorse,
dall'accumulo di stato interno anomalo o da derive lente nei pattern di traffico, non attiva
alcuna procedura di ricertificazione negli schemi tradizionali, poiché il codice rimane
invariato e il certificato non è ancora formalmente scaduto. Il risultato è una finestra di
invalidità silenziosa del certificato durante la quale il documento continua ad attestare la
conformità di un sistema che, nella pratica, non soddisfa più le proprietà per cui è stato
certificato.

Una terza limitazione è l'approccio alla selezione delle evidenze. Gli schemi tradizionali 
tendono infatti a focalizzare la valutazione esclusivamente su artefatti software, 
come il codice eseguibile e le configurazioni statiche. Così facendo, vengono spesso trascurati 
elementi che incidono in modo rilevante sulla qualità delle proprietà certificate; 
ad esempio, il processo di sviluppo adottato o le procedure di verifica applicate. 
Il risultato è la produzione di certificati che riflettono solo una parte dell'informazione 
necessaria a supportare decisioni accurate sulla compliance di un sistema distribuito~\cite{anisetti2023multidim}.

\subsection{Certificazione Continua e Multidimensionale}
\label{subsec:cert_avanzata}

La ricerca ha risposto ai limiti dei modelli tradizionali attraverso due direzioni complementari
che ampliano sia la dinamica temporale del processo certificativo sia l'ampiezza delle
dimensioni considerate nella valutazione.

La \textbf{certificazione continua} estende il processo di certificazione in un ciclo
\textbf{\gls{MAPE}} che monitora continuamente il comportamento del sistema comparandolo con
il modello di comportamento normale incluso nel certificato corrente. In questo schema, il
certificato $C_t$ al tempo $t$ integra esplicitamente un \emph{system model} $B$, ovvero un
modello adattivo del comportamento normale del sistema, costruito tramite tecniche formali, 
come Petri net o Abstract State Machine, oppure tecniche di machine learning applicate ai dati osservabili, a
seconda del contesto applicativo, oltre a un riferimento al certificato precedente
$C_{t-1}$~\cite{anisetti2023continuous}. In questo contesto, la ricertificazione è innescata
non da eventi discreti predefiniti, ma dalla rilevazione di una variazione comportamentale
misurata come scostamento tra il comportamento osservato e il modello $B$. Tale meccanismo
permette di ridurre i falsi positivi tipici dei trigger basati su semplici modifiche del codice
e consente di intercettare derive comportamentali che gli schemi convenzionali non riuscirebbero
a individuare.

La \textbf{certificazione multidimensionale}, invece, non si limita a valutare i soli artefatti software, 
bensì considera anche dimensioni aggiuntive, che influenzano la qualità delle proprietà certificate~\cite{anisetti2023multidim}. 
Una dimensione $D$ raggruppa attributi la cui inclusione nella valutazione, oltre agli attributi 
degli artefatti eseguibili, migliora l'accuratezza complessiva della valutazione di conformità. La certificazione, 
in questo caso, si basa sulla composizione dei risultati ottenuti nelle singole dimensioni. 
Così facendo è possibile ottenere un quadro più accurato della conformità attesa. 
L'approccio supporta inoltre la selezione dei servizi sulla base di più criteri di certificazione, 
rendendo così le decisioni più solide in contesti operativi reali.

% =========================================================
\section{Scenario di Riferimento: Sistema di Telemedicina Critica}
\label{sec:scenario_riferimento}
% =========================================================

Il formalismo introdotto nel Capitolo~\ref{chap:metodologia} richiede un sistema applicativo concreto su cui
esemplificare i costrutti del framework. La scelta è caduta su un sistema di telemedicina critica 
in quanto combina due caratteristiche necessarie a questa funzione: requisiti di compliance con vincoli 
quantitativi assoluti, e una topologia in cui due flussi applicativi con livelli di criticità diversi 
condividono risorse di backend. Questa seconda caratteristica è ciò che rende lo scenario importante 
per la dimostrazione del layer semantico; difatti, in un sistema del genere, un problema nel flusso meno 
critico può minacciare le garanzie del flusso più critico attraverso una risorsa condivisa.

\subsection{Architettura del Sistema e Flussi di Elaborazione}
\label{subsec:architettura_telemedicina}

L'architettura è costituita da cinque microservizi con funzionalità distinte:

\begin{itemize}
  \item \textbf{API Gateway (SGW)}: funge da punto di ingresso per le richieste esterne; 
      gestisce autenticazione, autorizzazione e instradamento verso i sottosistemi;
  \item \textbf{Data Ingestion Service (SIN)}: riceve i flussi di dati biomedici dai 
      dispositivi di monitoraggio remoto, li normalizza nel formato interno e li indirizza verso la persistenza; 
  \item \textbf{Database Service (SDB)}: rappresenta una risorsa di persistenza condivisa, 
      implementata con un pool di connessioni gestito dalla \gls{JVM};
  \item \textbf{Alert Service (SAL)}: recupera i dati da SDB, applica regole cliniche 
      predefinite e genera notifiche verso il personale sanitario;
  \item \textbf{Billing Service (SBL)}: gestisce fatturazione, archiviazione documentale 
      e procedure amministrative.
\end{itemize}

I cinque microservizi sono connessi da due flussi con criticità distinte. 
Il \textbf{flusso clinico} attraversa in sequenza i servizi SGW, SIN, SDB e SAL: 
ogni dato biomedico ricevuto dal gateway percorre l'intera catena fino alla 
generazione della notifica clinica. Il \textbf{flusso amministrativo} collega 
invece SGW al servizio di fatturazione SBL, che a sua volta accede a SDB per 
le operazioni di scrittura. La condivisione di SDB tra i due flussi costituisce 
il punto di interferenza strutturale del sistema, analizzato nella 
Sezione~\ref{subsec:compliance_sets}. La topologia complessiva è illustrata 
nella Figura~\ref{fig:telemedicina_grafo}.

\begin{figure}[H]
	\centering
	\includegraphics[width = 1.0\textwidth, trim={0 0 0 15mm}, clip]{images/telemedicina_grafo}
	\caption[Architettura del sistema di telemedicina]{Architettura del sistema di telemedicina}
	\label{fig:telemedicina_grafo}
\end{figure}

\subsection{Metriche di Osservazione del Sistema}
\label{subsec:metriche_osservazione}

Il sistema è implementato su runtime Java con framework Spring Boot, scelta che determina la
disponibilità di metriche di stato interno specifiche della \gls{JVM}. Le metriche di nodo
$\mathbf{x}_v(t) = [\text{cpu}_v(t),\, \text{mem}_v(t),\, \text{pool}_v(t),\, \text{gc}_v(t)]$
permettono di monitorare quattro dimensioni fondamentali di ciascun microservizio: utilizzo
della \gls{CPU}, occupazione dell'heap \gls{JVM}, saturazione del connection pool 
verso il database (\texttt{pool\_percent}), e la durata delle pause del garbage collector (\texttt{gc\_pause\_ms}).\\
Mentre l'utilizzo di \gls{CPU} e memoria rappresentano uno standard in qualsiasi sistema di
monitoraggio moderno, le restanti due metriche sono strettamente legate al runtime Java e
rivelano condizioni di degrado latente. Ad esempio, un nodo
può avere il connection pool saturo all'85\% senza che la latenza degli archi incidenti abbia
ancora violato le proprie soglie; in questo caso, le richieste vengono ancora servite, ma con tempi
di attesa crescenti nelle code di connessione. Queste metriche interne consentono di anticipare
la violazione rispetto al momento in cui essa diventa visibile sugli archi. \\
Le metriche di arco $\mathbf{x}_e(t) = [L_e(t),\, \varepsilon_e(t),\, \theta_e(t)]$
descrivono invece la qualità di ogni singola interazione attraverso la latenza media delle chiamate $L_e$, 
il tasso di errore $\varepsilon_e$ inteso come frazione di richieste fallite, e il throughput $\theta_e$ 
misurato in richieste al secondo. La Tabella~\ref{tab:metriche_telemedicina} riassume la classificazione completa.

\begin{table}[htbp]
\centering
\caption{Metriche del sistema di telemedicina.}
\label{tab:metriche_telemedicina}
\begin{tabular}{llll}
\toprule
\textbf{Simbolo} & \textbf{Descrizione} & \textbf{Unità} & \textbf{Natura} \\
\midrule
$\text{cpu}_v(t)$  & Utilizzo \gls{CPU}           & \%    & Stato (nodo) \\
$\text{mem}_v(t)$  & Heap \gls{JVM} occupato       & MB    & Stato (nodo) \\
$\text{pool}_v(t)$ & Saturazione connection pool   & \%    & Stato (nodo) \\
$\text{gc}_v(t)$   & Pause garbage collector       & ms    & Stato (nodo) \\
$L_e(t)$           & Latenza media arco            & ms    & Interazione (arco) \\
$\varepsilon_e(t)$ & Tasso di errore               & ratio & Interazione (arco) \\
$\theta_e(t)$      & Throughput                    & req/s & Interazione (arco) \\
\bottomrule
\end{tabular}
\end{table}

\subsection{Proprietà Certificate e Nodi Condivisi}
\label{subsec:compliance_sets}

Il sistema deve garantire il rispetto di due proprietà non funzionali certificate,
caratterizzate da requisiti distinti. La prima proprietà, denominata \textbf{Real-time
Dependability}, impone che la latenza end-to-end del flusso clinico non superi i 500\,ms e
che la disponibilità del percorso superi il 99,99\% su base mensile. Il soddisfacimento di
tale proprietà dipende dalla conformità congiunta dei servizi SGW, SIN, SDB e SAL; difatti, 
un malfunzionamento di uno qualunque di essi violerebbe la proprietà sull'intero percorso.
La seconda proprietà, denominata \textbf{Billing Compliance}, richiede
la corretta esecuzione di ogni operazione di fatturazione e la coerenza della base documentale.
I componenti coinvolti sono SGW, SBL e SDB.

Il confronto tra i due insiemi rivela immediatamente che SGW e SDB appartengono a entrambe le proprietà. 
In particolare, SDB rappresenta il nodo di interferenza strutturale critico tra i due flussi. 
Un incremento sostenuto del carico amministrativo sul flusso di fatturazione può infatti saturare 
il \emph{connection pool} del database, rallentando l'accesso di SIN alla persistenza. 
Si genera così una minaccia per la proprietà \textbf{Real-time Dependability} attraverso una catena 
causale che ha origine interamente al di fuori del percorso clinico. 
Questa interferenza non è identificabile analizzando i flussi in isolamento; essa emerge invece come 
informazione topologica dal confronto tra i due insiemi di componenti.
Il valore architetturale del layer semantico risiede proprio nella capacità di rendere
esplicita, già in fase di progettazione, un'interferenza che i sistemi di assessment
tradizionali potrebbero rilevare solo empiricamente. In certi contesti, un rilevamento tardivo
renderebbe inutile qualsiasi tentativo di intervento. Questo scenario verrà utilizzato come esempio illustrativo nel Capitolo~\ref{chap:metodologia} per mostrare
come la pipeline predittiva operi in presenza di un'interferenza cross-property.